% Options for packages loaded elsewhere
% Options for packages loaded elsewhere
\PassOptionsToPackage{unicode}{hyperref}
\PassOptionsToPackage{hyphens}{url}
\PassOptionsToPackage{dvipsnames,svgnames,x11names}{xcolor}
%
\documentclass[
  japanese,
  lualatex,
  ja=standard,
  12pt,
  a4paper]{ltjsarticle}
\usepackage{xcolor}
\usepackage{amsmath,amssymb}
\setcounter{secnumdepth}{5}
\usepackage{iftex}
\ifPDFTeX
  \usepackage[T1]{fontenc}
  \usepackage[utf8]{inputenc}
  \usepackage{textcomp} % provide euro and other symbols
\else % if luatex or xetex
  \usepackage{unicode-math} % this also loads fontspec
  \defaultfontfeatures{Scale=MatchLowercase}
  \defaultfontfeatures[\rmfamily]{Ligatures=TeX,Scale=1}
\fi
\usepackage{lmodern}
\ifPDFTeX\else
  % xetex/luatex font selection
\fi
% Use upquote if available, for straight quotes in verbatim environments
\IfFileExists{upquote.sty}{\usepackage{upquote}}{}
\IfFileExists{microtype.sty}{% use microtype if available
  \usepackage[]{microtype}
  \UseMicrotypeSet[protrusion]{basicmath} % disable protrusion for tt fonts
}{}
\makeatletter
\@ifundefined{KOMAClassName}{% if non-KOMA class
  \IfFileExists{parskip.sty}{%
    \usepackage{parskip}
  }{% else
    \setlength{\parindent}{0pt}
    \setlength{\parskip}{6pt plus 2pt minus 1pt}}
}{% if KOMA class
  \KOMAoptions{parskip=half}}
\makeatother
% Make \paragraph and \subparagraph free-standing
\makeatletter
\ifx\paragraph\undefined\else
  \let\oldparagraph\paragraph
  \renewcommand{\paragraph}{
    \@ifstar
      \xxxParagraphStar
      \xxxParagraphNoStar
  }
  \newcommand{\xxxParagraphStar}[1]{\oldparagraph*{#1}\mbox{}}
  \newcommand{\xxxParagraphNoStar}[1]{\oldparagraph{#1}\mbox{}}
\fi
\ifx\subparagraph\undefined\else
  \let\oldsubparagraph\subparagraph
  \renewcommand{\subparagraph}{
    \@ifstar
      \xxxSubParagraphStar
      \xxxSubParagraphNoStar
  }
  \newcommand{\xxxSubParagraphStar}[1]{\oldsubparagraph*{#1}\mbox{}}
  \newcommand{\xxxSubParagraphNoStar}[1]{\oldsubparagraph{#1}\mbox{}}
\fi
\makeatother


\usepackage{longtable,booktabs,array}
\newcounter{none} % for unnumbered tables
\usepackage{calc} % for calculating minipage widths
% Correct order of tables after \paragraph or \subparagraph
\usepackage{etoolbox}
\makeatletter
\patchcmd\longtable{\par}{\if@noskipsec\mbox{}\fi\par}{}{}
\makeatother
% Allow footnotes in longtable head/foot
\IfFileExists{footnotehyper.sty}{\usepackage{footnotehyper}}{\usepackage{footnote}}
\makesavenoteenv{longtable}
\usepackage{graphicx}
\makeatletter
\newsavebox\pandoc@box
\newcommand*\pandocbounded[1]{% scales image to fit in text height/width
  \sbox\pandoc@box{#1}%
  \Gscale@div\@tempa{\textheight}{\dimexpr\ht\pandoc@box+\dp\pandoc@box\relax}%
  \Gscale@div\@tempb{\linewidth}{\wd\pandoc@box}%
  \ifdim\@tempb\p@<\@tempa\p@\let\@tempa\@tempb\fi% select the smaller of both
  \ifdim\@tempa\p@<\p@\scalebox{\@tempa}{\usebox\pandoc@box}%
  \else\usebox{\pandoc@box}%
  \fi%
}
% Set default figure placement to htbp
\def\fps@figure{htbp}
\makeatother



\ifLuaTeX
\usepackage[bidi=basic,shorthands=off,provide=*]{babel}
\else
\usepackage[bidi=default,shorthands=off,provide=*]{babel}
\fi
\ifLuaTeX
  \usepackage{selnolig} % disable illegal ligatures
\fi


\setlength{\emergencystretch}{3em} % prevent overfull lines

\providecommand{\tightlist}{%
  \setlength{\itemsep}{0pt}\setlength{\parskip}{0pt}}



 


\usepackage{needspace}
\makeatletter
\@ifpackageloaded{tcolorbox}{}{\usepackage[skins,breakable]{tcolorbox}}
\@ifpackageloaded{fontawesome5}{}{\usepackage{fontawesome5}}
\definecolor{quarto-callout-color}{HTML}{909090}
\definecolor{quarto-callout-note-color}{HTML}{0758E5}
\definecolor{quarto-callout-important-color}{HTML}{CC1914}
\definecolor{quarto-callout-warning-color}{HTML}{EB9113}
\definecolor{quarto-callout-tip-color}{HTML}{00A047}
\definecolor{quarto-callout-caution-color}{HTML}{FC5300}
\definecolor{quarto-callout-color-frame}{HTML}{acacac}
\definecolor{quarto-callout-note-color-frame}{HTML}{4582ec}
\definecolor{quarto-callout-important-color-frame}{HTML}{d9534f}
\definecolor{quarto-callout-warning-color-frame}{HTML}{f0ad4e}
\definecolor{quarto-callout-tip-color-frame}{HTML}{02b875}
\definecolor{quarto-callout-caution-color-frame}{HTML}{fd7e14}
\makeatother
\makeatletter
\@ifpackageloaded{caption}{}{\usepackage{caption}}
\AtBeginDocument{%
\ifdefined\contentsname
  \renewcommand*\contentsname{目次}
\else
  \newcommand\contentsname{目次}
\fi
\ifdefined\listfigurename
  \renewcommand*\listfigurename{図一覧}
\else
  \newcommand\listfigurename{図一覧}
\fi
\ifdefined\listtablename
  \renewcommand*\listtablename{表一覧}
\else
  \newcommand\listtablename{表一覧}
\fi
\ifdefined\figurename
  \renewcommand*\figurename{図}
\else
  \newcommand\figurename{図}
\fi
\ifdefined\tablename
  \renewcommand*\tablename{表}
\else
  \newcommand\tablename{表}
\fi
}
\@ifpackageloaded{float}{}{\usepackage{float}}
\floatstyle{ruled}
\@ifundefined{c@chapter}{\newfloat{codelisting}{h}{lop}}{\newfloat{codelisting}{h}{lop}[chapter]}
\floatname{codelisting}{コード}
\newcommand*\listoflistings{\listof{codelisting}{コード一覧}}
\makeatother
\makeatletter
\makeatother
\makeatletter
\@ifpackageloaded{caption}{}{\usepackage{caption}}
\@ifpackageloaded{subcaption}{}{\usepackage{subcaption}}
\makeatother
\usepackage{bookmark}
\IfFileExists{xurl.sty}{\usepackage{xurl}}{} % add URL line breaks if available
\urlstyle{same}
% fallback for those not using the hyperref driver hyperxmp:
\makeatletter
\@ifundefined{xmpquote}{\newcommand{\xmpquote}[1]{#1}}{}
\makeatother
\hypersetup{
  pdftitle={第5回：政府部門と財政政策},
  pdfauthor={宮﨑憲治},
  pdflang={ja},
  colorlinks=true,
  linkcolor={blue},
  filecolor={Maroon},
  citecolor={Blue},
  urlcolor={Blue},
  pdfcreator={LaTeX via pandoc}}


\title{第5回：政府部門と財政政策}
\usepackage{etoolbox}
\makeatletter
\providecommand{\subtitle}[1]{% add subtitle to \maketitle
  \apptocmd{\@title}{\par {\large #1 \par}}{}{}
}
\makeatother
\subtitle{マクロ経済学B}
\author{宮﨑憲治}
\date{2026年度}
\begin{document}
\maketitle

\renewcommand*\contentsname{目次}
{
\setcounter{tocdepth}{3}
\tableofcontents
}

\section{講義の目的}\label{ux8b1bux7fa9ux306eux76eeux7684}

\subsection{本講義で探究する中心的な問い}\label{ux672cux8b1bux7fa9ux3067ux63a2ux7a76ux3059ux308bux4e2dux5fc3ux7684ux306aux554fux3044}

本講義では、第4回で構築した資本蓄積のない柔軟価格（Flexible
Price）モデルに「\textbf{政府部門（Government
Sector）}」を導入し、財政政策がマクロ経済に与える影響を理論的に解明します。

マクロ経済学において政府の役割を分析する際、最も重要な論点は「\textbf{政府による実物資源の直接的な吸収（政府購入
\(G_t\)）}」と「\textbf{民間主体間における純粋な購買力の移転（租税
\(TX_t\) や移転支出
\(TR_t\)）}」を明確に峻別することです。本講義では、国債発行を通じた財源調達のタイミングの選択が実物配分に影響を与えないという「\textbf{バロー・リカードの中立命題（Ricardian
Equivalence）}」の論理構造を確認した上で、外生的な政府支出ショックが民間消費や労働供給、総産出に及ぼす波及メカニズムを対数線形近似を用いて詳細に解明します。

さらに講義の後半では、一括税とは異なり人々の限界的な意思決定に歪みをもたらす「\textbf{比例的労働所得税}」を取り上げ、税率の引き上げが労働意欲や課税ベースを縮小させるメカニズムを通じて、税率と税収の間に逆U字型の関係が生じる「\textbf{ラッファーカーブ（Laffer
Curve）}」を静学的に導出・考察します。

本講義の具体的な到達目標は、以下の4点です。

\begin{enumerate}
\def\labelenumi{\arabic{enumi}.}
\tightlist
\item
  \textbf{政府予算制約と公債動学の理解}：単年度のフロー予算制約式からプライマリーバランスを定義し、家計の消費オイラー方程式と横断性条件を用いて政府の通時的予算制約式を導出できるようになる。
\item
  \textbf{財源調達と資源吸収の区別}：一括税・一括移転と政府購入（政府支出）がマクロ経済の実物配分に与える影響の相違を、リカードの中立命題の成立要件とともに論理的に説明できるようになる。
\item
  \textbf{政府支出ショックの対数線形解析}：資本蓄積のない固定資本モデルにおいて政府支出を含む均衡体系を対数線形化し、政府購入乗数
  \(\zeta_g\) および各マクロ変数の反応係数を導出できるようになる。
\item
  \textbf{税制のくさびとラッファーカーブの導出}：労働所得税と消費税が労働市場に生み出す「くさび」や税の帰着の等価性を理解し、労働所得税率と税収の関係を表すラッファーカーブおよび最大税収率を解析的に導出できるようになる。
\end{enumerate}

本講義では、以下の表に示すように、大きく2つの政策実験を明確に区別して進めます。

{\def\LTcaptype{none} % do not increment counter
\begin{longtable}[]{@{}
  >{\raggedright\arraybackslash}p{(\linewidth - 4\tabcolsep) * \real{0.3333}}
  >{\raggedright\arraybackslash}p{(\linewidth - 4\tabcolsep) * \real{0.3333}}
  >{\raggedright\arraybackslash}p{(\linewidth - 4\tabcolsep) * \real{0.3333}}@{}}
\toprule\noalign{}
\begin{minipage}[b]{\linewidth}\raggedright
政策実験
\end{minipage} & \begin{minipage}[b]{\linewidth}\raggedright
政策設定の前提
\end{minipage} & \begin{minipage}[b]{\linewidth}\raggedright
主たる分析対象と目的
\end{minipage} \\
\midrule\noalign{}
\endhead
\bottomrule\noalign{}
\endlastfoot
\textbf{政府購入ショックの動学分析} & 政府購入 \(G_t\)
が確率的に変動し、一括税または国債で調達 &
民間消費、労働投入、産出量、実質賃金への波及効果と乗数 \\
\textbf{労働所得税の静学実験} & \(G=0, B=0\) のもとで、税収を一括移転
\(TR=TX\) で還付 &
比例税率と課税ベースの縮小、最大税収率（ラッファーカーブ） \\
\end{longtable}
}

図~\ref{fig-lecture05-overview}
は、固定資本モデルの実物体系を基礎として、政府の予算制約、政府支出ショックの伝播、そして税制の帰着とラッファーカーブへと展開していく本講義全体の論理的構成を示したものです。

\begin{figure}[H]

\centering{

\includegraphics[width=1\linewidth,height=\textheight,keepaspectratio]{lecture05_files/mediabag/figures/lecture05/lecture05_overview_infographic.pdf}

}

\caption{\label{fig-lecture05-overview}第5回講義の全体像：固定資本の実物経済モデルをベースに、政府予算制約の導入、政府支出ショックの波及効果、税の帰着とラッファーカーブを体系的に分析する流れ。}

\end{figure}%

\section{表記と基本設定}\label{ux8868ux8a18ux3068ux57faux672cux8a2dux5b9a}

本講義ノートでは、慣例に従い、変数の\textbf{水準値を大文字}で、定常状態からの\textbf{対数偏差（Log-deviation）を小文字}で表記します。ただし、政府購入
\(G_t\)
の変動については、定常状態における政府支出比率がゼロ（\(G=0\)）の場合でも対数近似が破綻せず安定して扱えるよう、定常状態の総産出
\(Y\) で割った\textbf{産出比率としての乖離} \[
g_t \equiv \frac{G_t - G}{Y}
\] として定義します。

本ノートで使用する主要な水準変数の記号一覧は以下の通りです。

{\def\LTcaptype{none} % do not increment counter
\begin{longtable}[]{@{}
  >{\raggedright\arraybackslash}p{(\linewidth - 6\tabcolsep) * \real{0.2500}}
  >{\raggedright\arraybackslash}p{(\linewidth - 6\tabcolsep) * \real{0.2500}}
  >{\raggedright\arraybackslash}p{(\linewidth - 6\tabcolsep) * \real{0.2500}}
  >{\raggedright\arraybackslash}p{(\linewidth - 6\tabcolsep) * \real{0.2500}}@{}}
\toprule\noalign{}
\begin{minipage}[b]{\linewidth}\raggedright
記号
\end{minipage} & \begin{minipage}[b]{\linewidth}\raggedright
経済学的意味
\end{minipage} & \begin{minipage}[b]{\linewidth}\raggedright
記号
\end{minipage} & \begin{minipage}[b]{\linewidth}\raggedright
経済学的意味
\end{minipage} \\
\midrule\noalign{}
\endhead
\bottomrule\noalign{}
\endlastfoot
\(C_t\) & 民間実質消費水準 & \(N_t\) & 労働投入量（就業時間） \\
\(Y_t\) & 総産出水準（実質GDP） & \(G_t\) &
実質政府購入（政府支出水準） \\
\(W_t\) & 実質賃金率 & \(\bar{K}=1\) &
企業が物理的に保有する固定資本ストック（\(1\) に正規化） \\
\(B_t(h)\) & 家計 \(h\) の期末実質債券保有残高 & \(B_t\) &
経済全体における期末実質政府債務（国債発行残高） \\
\(\Theta_t(h)\) & 家計 \(h\) の企業株式（持分）保有株数 & \(V_t\) &
配当落ち後の企業株式時価総額（株価） \\
\(D_t\) & 企業が支払う実質利潤（実質配当） & \(TX_t\) &
政府が徴収する実質総税収 \\
\(TR_t\) & 政府が給付する実質一括移転 & \(PB_t\) &
プライマリーバランス（基礎的財政収支） \\
\(R_t^N\) & 1期間名目総利子率（グロス名目利子率） & \(\Pi_t\) &
粗インフレ率（\(P_t/P_{t-1}\)） \\
\(P_t\) & 財市場における絶対物価水準 & \(\mathbb{E}_t\) & 時点 \(t\)
の情報集合に基づく条件付き期待値演算子 \\
\(a_t\) & 全要素生産性（技術水準の対数値） & \(e_t^a, e_t^g\) &
技術ショックおよび政府支出ショックの白色雑音イノベーション \\
\(C, N, Y, G, W\) & 対応する各変数の確定的定常状態における水準値 & & \\
\end{longtable}
}

また、経済の構造を規定する基礎パラメータは以下の通りです。

{\def\LTcaptype{none} % do not increment counter
\begin{longtable}[]{@{}
  >{\raggedright\arraybackslash}p{(\linewidth - 2\tabcolsep) * \real{0.5000}}
  >{\raggedright\arraybackslash}p{(\linewidth - 2\tabcolsep) * \real{0.5000}}@{}}
\toprule\noalign{}
\begin{minipage}[b]{\linewidth}\raggedright
記号
\end{minipage} & \begin{minipage}[b]{\linewidth}\raggedright
定義・経済学的意味
\end{minipage} \\
\midrule\noalign{}
\endhead
\bottomrule\noalign{}
\endlastfoot
\(0 < \beta < 1\) & 家計の主観的割引因子 \\
\(\gamma > 0\) &
相対的リスク回避度（異時点間消費代替弾力性の逆数）。同弾力性は
\(1/\gamma\) \\
\(\varphi > 0\) &
労働の限界不効用パラメータ（Frisch労働供給弾力性の逆数）。同弾力性は
\(1/\varphi\) \\
\(0 \leq \alpha < 1\) &
Cobb-Douglas生産関数における固定資本の分配率（資本シェア） \\
\(\nu = 1\) & 労働不効用の尺度パラメータ（計算の簡素化のため \(1\)
に正規化） \\
\(0 \leq G_Y \equiv G/Y < 1\) &
定常状態における政府支出の対産出比率（政府支出シェア） \\
\(|\rho_a| < 1, |\rho_g| < 1\) &
技術ショックおよび政府支出ショックの自己相関（持続性）パラメータ \\
\(\tau_N, \tau_C\) & 労働所得税率および消費税率（政策パラメータ） \\
\end{longtable}
}

さらに、本講義の展開において頻出する重要な派生変数は以下の通りです。なお、ショックに対する反応係数
\(\zeta_a, \zeta_g\)
やラッファーカーブに関連する係数については、それぞれの導出節において明示的に定義します。

{\def\LTcaptype{none} % do not increment counter
\begin{longtable}[]{@{}
  >{\raggedright\arraybackslash}p{(\linewidth - 4\tabcolsep) * \real{0.3333}}
  >{\raggedright\arraybackslash}p{(\linewidth - 4\tabcolsep) * \real{0.3333}}
  >{\raggedright\arraybackslash}p{(\linewidth - 4\tabcolsep) * \real{0.3333}}@{}}
\toprule\noalign{}
\begin{minipage}[b]{\linewidth}\raggedright
記号
\end{minipage} & \begin{minipage}[b]{\linewidth}\raggedright
定義式
\end{minipage} & \begin{minipage}[b]{\linewidth}\raggedright
経済学的意味
\end{minipage} \\
\midrule\noalign{}
\endhead
\bottomrule\noalign{}
\endlastfoot
\(g_t\) & \((G_t - G)/Y\) &
定常状態の産出水準で規格化した政府支出ショックの乖離幅 \\
\(Q_{t,t+i}\) & \(\beta^i U_{C,t+i}/U_{C,t}\) & 時点 \(t\) から時点
\(t+i\) までの多期間確率的割引因子（SDF） \\
\end{longtable}
}

\section{政府予算制約}\label{ux653fux5e9cux4e88ux7b97ux5236ux7d04}

政府は各期において、財市場から財を直接買い入れる実質政府購入 \(G_t\)
および民間家計への実質一括移転 \(TR_t\)
を支弁します。これらの歳出を賄う財源は、実質租税収入 \(TX_t\)
と新規の国債発行です。

\(B_t\) は、政府が時点 \(t\)
において新規に発行する実質国債の発行額（実質調達額）を表します。この国債は名目額面で発行され、1期間後の時点
\(t+1\) において名目総利子率 \(R_t^N\) を付して償還されるため、時点
\(t+1\) における事後的な実質返済額は \(R_t^N B_t / \Pi_{t+1}\)
となります。

したがって、各期 \(t\)
における政府の\textbf{単年度実質フロー予算制約式}は、次のように記述されます。
\begin{equation}\protect\phantomsection\label{eq-government-flow-budget}{
G_t + TR_t + \frac{R_{t-1}^N}{\Pi_t} B_{t-1} = TX_t + B_t.
}\end{equation}
左辺は政府の当期総支出（実質政府購入、一括移転、および前期発行国債の実質元利返済額）であり、右辺は利用可能な総実質収入（当期税収および新規国債発行額）を表しています。

ここで、国債の利払い・償還費を除く純粋な政策的収支である\textbf{プライマリーバランス（Primary
Balance: 基礎的財政収支）}を \[
PB_t \equiv TX_t - TR_t - G_t
\]
と定義すると、政府のフロー予算制約式は次のような\textbf{公債動学方程式}として書き直すことができます。
\begin{equation}\protect\phantomsection\label{eq-government-debt-law}{
B_t = \frac{R_{t-1}^N}{\Pi_t} B_{t-1} - PB_t.
}\end{equation} この式から明らかなように、当期末の政府債務残高 \(B_t\)
が増加する要因は、前期からの利払いを含む債務負担（第1項）が大きいか、あるいはプライマリーバランスが赤字（\(PB_t < 0\)）であることのいずれか（または双方）です。

次に、政府の長期的な財政の持続可能性を特徴づける\textbf{通時的予算制約式（Intertemporal
Budget
Constraint）}を導出します。家計の無リスク名目債券に関するオイラー方程式から、確率的割引因子
\(Q_{t,t+1}\) を用いて \[
1 = \mathbb{E}_t \left[ Q_{t,t+1} \frac{R_t^N}{\Pi_{t+1}} \right]
\] が常に成立します。1期先（時点 \(t+1\)）における公債動学方程式を \[
B_{t+1} + PB_{t+1} = \frac{R_t^N}{\Pi_{t+1}} B_t
\] と書き改め、両辺に \(Q_{t,t+1}\) を掛けて時点 \(t\)
の情報集合に基づく条件付き期待値をとります。時点 \(t\) において \(B_t\)
はすでに確定した既知の値であるため、期待値の外に取り出すことができ、債券のオイラー方程式を適用すると
\[
\mathbb{E}_t \left[ Q_{t,t+1} (B_{t+1} + PB_{t+1}) \right] = B_t \mathbb{E}_t \left[ Q_{t,t+1} \frac{R_t^N}{\Pi_{t+1}} \right] = B_t
\] という極めて簡潔な関係が得られます。

この関係式を将来に向かって前向きに逐次反復代入し、「反復期待値の法則（Law
of Iterated Expectations）」を適用すると、任意の有限期間 \(N\)
に対して次式が成立します。 \[
B_t = \mathbb{E}_t \left[ \sum_{i=1}^N Q_{t,t+i} PB_{t+i} + Q_{t,t+N} B_{t+N} \right].
\]
ここで、政府が将来にわたって債務を無限に借り換え続けて元利返済を永遠に先送りすることを排除するため、以下の\textbf{公債に関する横断性条件（Transversality
Condition: TVC）}を課します。
\begin{equation}\protect\phantomsection\label{eq-government-transversality}{
\lim_{i \to \infty} \mathbb{E}_t \left[ Q_{t,t+i} B_{t+i} \right] = 0.
}\end{equation} この条件が満たされるとき、\(N \to \infty\)
の極限をとることで、政府の\textbf{通時的予算制約式}が得られます。
\begin{equation}\protect\phantomsection\label{eq-government-intertemporal-budget}{
B_t = \mathbb{E}_t \sum_{i=1}^\infty Q_{t,t+i} PB_{t+i}.
}\end{equation}
この方程式は、「\textbf{現在残存している政府債務の実質残高 \(B_t\)
は、将来にわたって生み出されるプライマリーバランス黒字流列の確率的割引現在価値の総和に厳密に等しくなければならない}」という財政の長期均衡条件を意味しています。

ただし留意すべき点として、この通時的予算制約式単体から直ちに「国債発行が経済に中立である」という結論が導かれるわけではありません。家計側の通時的予算制約および資産市場の均衡と統合し、「将来の一括税の引き上げによって政府債務が完全に償還される」という財政政策の完結条件（財政政策の閉じ方）が満たされて初めて、公債発行の中立性が成立します。

\section{家計・企業・均衡}\label{ux5bb6ux8a08ux4f01ux696dux5747ux8861}

\subsection{家計の行動}\label{ux5bb6ux8a08ux306eux884cux52d5}

家計のミクロ的最適化行動は、基本的に第4回と同様です。経済全体の家計の総測度および発行株式総数を
\(1\) に正規化し、各家計に一括税 \(TX_t\) および一括移転 \(TR_t\)
が課される状況を考えます。

代表的家計 \(h\)
の各期における実質フロー予算制約式は、以下のように記述されます。 \[
\begin{aligned}
&C_t(h) + B_t(h) + \Theta_t(h) V_t + TX_t\\
&= \frac{R_{t-1}^N}{\Pi_t} B_{t-1}(h)
 + \Theta_{t-1}(h) (V_t + D_t)\\
&\quad + W_t N_t(h) + TR_t.
\end{aligned}
\]
左辺は消費、新規の国債購入、新規の株式購入、および一括税の支払いからなる総支出であり、右辺は前期保有国債の実質元利受取額、前期保有株式の売却評価額および受取配当、労働所得、ならびに政府からの一括移転からなる総実質資源です。家計の総測度が
\(1\) であるため、各家計に一律に課される \(TX_t\) と \(TR_t\)
は、そのままマクロ集計値と一致します。

すべての家計が同一の初期資産を持ち、同一の選択を行う対称均衡のもとで、家計の効用最大化問題の一階条件（FOC）は以下の3式となります。
\[
\begin{aligned}
W_t &= -\frac{U_{N,t}}{U_{C,t}},\\
1 &= \mathbb{E}_t \left[ Q_{t,t+1} \frac{R_t^N}{\Pi_{t+1}} \right],\\
V_t &= \mathbb{E}_t \left[ Q_{t,t+1} (V_{t+1} + D_{t+1}) \right].
\end{aligned}
\] ここで注目すべき重要な事実は、\textbf{一括税 \(TX_t\) および一括移転
\(TR_t\)
は、これらの一階条件（限界代替率やオイラー方程式）に直接現れない}という点です。これは、一括税・移転が家計の限界的な労働意欲や貯蓄インセンティブを歪めないことを示しています。ただし、これは一括税がマクロ経済に何の影響も与えないという意味ではありません。後述するように、政府と家計の予算制約が統合される結果、一括税は生涯可処分所得の変動を通じた「所得効果（富効果）」を通じてのみ実物配分に作用します。

\subsection{企業の行動}\label{ux4f01ux696dux306eux884cux52d5}

企業の利潤最大化行動も第4回と完全に同様です。代表的な完全競争企業は、固定された資本ストック
\(\bar{K}=1\) を保有し、毎期ごとに利潤を最大化するように労働投入量
\(N_t\) を決定します。

企業の生産関数および当期の実質利潤（実質配当）\(D_t\)
は、次式で与えられます。 \[
Y_t = \exp(a_t) F(\bar{K}, N_t) = \exp(a_t) F(1, N_t), \qquad D_t = Y_t - W_t N_t.
\]
プライス・テイカーとして行動する企業の一階条件は、労働の限界生産力（\(MPN_t\)）を市場の実質賃金率
\(W_t\) に等置することです。 \[
W_t = \exp(a_t) F_N(\bar{K}, N_t).
\]
したがって、家計の労働供給条件と等置することにより、労働市場の均衡条件
\[
-\frac{U_{N,t}}{U_{C,t}} = \exp(a_t) F_N(\bar{K}, N_t)
\] が得られます。

\subsection{政府支出を含む均衡と総資源制約式}\label{ux653fux5e9cux652fux51faux3092ux542bux3080ux5747ux8861ux3068ux7dcfux8cc7ux6e90ux5236ux7d04ux5f0f}

本講義で導入する政府購入 \(G_t\)
は、家計の効用を直接高めることもなく、また企業の生産性を向上させるインフラとしても機能しない「純粋に実物資源を消費・吸収する活動」として定式化されます。

金融市場における国債および株式の市場清算条件は、以下の通りです。 \[
\int_0^1 B_t(h)\,dh = B_t, \qquad \int_0^1 \Theta_t(h)\,dh = 1.
\]
代表的家計のフロー予算制約式を経済全体で集計すると、株式の売買代金は相殺され、企業の利潤分配関係
\(D_t + W_t N_t = Y_t\)
を代入することにより、集計された家計部門の収支関係 \[
C_t + B_t + TX_t = \frac{R_{t-1}^N}{\Pi_t} B_{t-1} + Y_t + TR_t
\]
が得られます。この式を政府の単年度フロー予算制約式（式~\ref{eq-government-flow-budget}）と統合すると、国債残高
\(B_t\)、税収 \(TX_t\)、および移転支出 \(TR_t\)
が左右両辺で完全に相殺されます。その結果、投資が存在しない閉鎖経済における財市場の需給均衡条件、すなわち\textbf{総資源制約式（Resource
Constraint）}が導かれます。
\begin{equation}\protect\phantomsection\label{eq-resource-constraint-government}{
Y_t = C_t + G_t.
}\end{equation} この式は、国内で生産された総実質付加価値 \(Y_t\)
が、民間消費 \(C_t\) と政府購入 \(G_t\)
の2つの用途に完全に配分されることを示しています。

定常状態における政府支出の対産出比率を \[
G_Y \equiv \frac{G}{Y}
\] と定義し、政府支出ショック \(g_t\) を定常産出 \(Y\)
に対する乖離として
\begin{equation}\protect\phantomsection\label{eq-government-spending-shock}{
g_t \equiv \frac{G_t - G}{Y}
}\end{equation} と定義します。外生的な政府支出ショック \(g_t\)
は、以下の標準的な1階自己回帰過程（AR(1)過程）に従うと仮定します。 \[
g_t = \rho_g g_{t-1} + e_t^g, \qquad |\rho_g| < 1, \qquad \mathbb{E}_{t-1}[e_t^g] = 0.
\]
以上より、本モデルにおける実物配分（\(W_t, N_t, Y_t, C_t, g_t\)）を決定する静学的均衡体系は、以下の5本の方程式に集約されます。
\[
\begin{aligned}
W_t &= -\frac{U_{N,t}}{U_{C,t}},\\
W_t &= \exp(a_t) F_N(\bar{K}, N_t),\\
Y_t &= \exp(a_t) F(\bar{K}, N_t),\\
Y_t &= C_t + G_t,\\
g_t &= \frac{G_t - G}{Y}.
\end{aligned}
\]
完全な動学的一般均衡体系を記述するためには、これら5本の実物側の方程式に加えて、家計の消費オイラー方程式、政府の予算制約式、資産市場の清算条件、および財政政策の完結条件を連立させる必要があります。

\begin{tcolorbox}[enhanced jigsaw, arc=.35mm, bottomrule=.15mm, bottomtitle=1mm, breakable, colback=white, colbacktitle=quarto-callout-note-color!10!white, colframe=quarto-callout-note-color-frame, coltitle=black, left=2mm, leftrule=.75mm, opacityback=0, opacitybacktitle=0.6, rightrule=.15mm, title=\textcolor{quarto-callout-note-color}{\faInfo}\hspace{0.5em}{バロー・リカードの中立命題を支える5つの前提条件}, titlerule=0mm, toprule=.15mm, toptitle=1mm]

現代マクロ経済学において「\textbf{リカードの中立命題（Ricardian
Equivalence）}」として知られる定理（Barro,
1974）は、以下の5つの前提条件が満たされるときに厳密に成立します：

\begin{enumerate}
\def\labelenumi{\arabic{enumi}.}
\tightlist
\item
  \textbf{無限の計画期間}：代表的家計が無限期間を生きるか、あるいは親族への遺産動機を通じて世代間で利他主義的に結ばれていること。
\item
  \textbf{完全な金融市場}：家計が借入制約（流動性制約）に直面せず、政府と同一の金利で自由に通時的な貸借を行えること。
\item
  \textbf{一括固定税の存在}：政府の課税が一括税（Lump-sum
  Tax）によって行われ、労働や貯蓄の限界的意思決定を歪めないこと。
\item
  \textbf{財政の持続可能性}：政府の横断性条件が満たされ、国債が将来のプライマリーバランス黒字によって裏づけられていること。
\item
  \textbf{所与の政府購入経路}：政府による実質資源の吸収経路
  \(\{G_t\}_{t=0}^\infty\) が所与として一定に保たれていること。
\end{enumerate}

これらの条件が満たされるとき、政府が当期の支出を増税ではなく国債発行によって賄ったとしても、家計は「将来いつか同額（現在価値評価）の増税が行われる」と完全に合理的に予見します。したがって、国債発行によって手元に残った資金は富の増加とはみなされず、将来の増税に備えて全額が貯蓄に回されます。その結果、\textbf{財源調達手段の選択（税か国債か）は、民間消費、投資、産出量などの実物配分に一切の影響を与えません}。

しかし、現実の経済においてこれらの仮定のいずれか（特に借入制約や世代交代、あるいは歪みのある比例税）が満たされない場合、リカードの中立命題は崩れ、国債発行のタイミングが実物経済に重大な影響を及ぼすことになります。また、将来の増税が保証されない非リカード的な財政政策のもとでは、物価水準が財政の動学によって決定される「物価水準の財政理論（FTPL:
Fiscal Theory of the Price Level）」（Leeper,
1991）の枠組みが必要となります。

\end{tcolorbox}

\section{関数形と定常状態}\label{ux95a2ux6570ux5f62ux3068ux5b9aux5e38ux72b6ux614b}

\subsection{効用関数と生産技術の特定化}\label{ux52b9ux7528ux95a2ux6570ux3068ux751fux7523ux6280ux8853ux306eux7279ux5b9aux5316}

分析をさらに具体化するため、第4回と同様にマクロ経済学の標準的な関数形を特定化します。

家計の消費に関する期間効用関数 \(u(C)\)
としてCRRA効用関数を、労働の不効用としてFrisch型を採用します。 \[
U(C_t, N_t) = \frac{C_t^{1-\gamma}-1}{1-\gamma} - \frac{N_t^{1+\varphi}}{1+\varphi}, \qquad \gamma > 0, \quad \varphi > 0.
\] 企業の生産技術は、固定資本 \(\bar{K}=1\)
のもとでのCobb-Douglas生産関数です。 \[
Y_t = \exp(a_t) \bar{K}^\alpha N_t^{1-\alpha} = \exp(a_t) N_t^{1-\alpha}, \qquad 0 \leq \alpha < 1.
\] 第4回で確認した通り、固定資本が \(1\)
に正規化されていても資本の限界収益はゼロにはならず、企業の労働需要条件
\(W_t N_t = (1-\alpha) Y_t\) より、生産物のうち \(\alpha\)
の割合が株主への実質配当（固定資本の収益）として分配されます（\(D_t = \alpha Y_t\)）。

この特定化のもとで、主観的限界代替率 \(MRS_t = N_t^\varphi C_t^\gamma\)
と労働の限界生産力 \(MPN_t = (1-\alpha)\exp(a_t)N_t^{-\alpha}\)
を等置すると、労働市場の均衡条件は次式で与えられます。 \[
N_t^\varphi C_t^\gamma = (1-\alpha) \exp(a_t) N_t^{-\alpha}.
\] ここに総資源制約式
\(C_t = Y_t - G_t = \exp(a_t) N_t^{1-\alpha} - G_t\)
を代入すると、政府購入 \(G_t\)
が実物資源を吸収することにより、所与の産出量のもとで民間が利用可能な消費
\(C_t\)
が直接的に圧迫されることがわかります。消費水準の低下は消費の限界効用を押し上げ、家計に対して労働供給を増やそうとする強力な「\textbf{負の富効果（負の所得効果）}」をもたらします。本モデルでは投資による資本蓄積が存在しないため、政府購入の拡大によるクラウディング・アウトは、民間投資ではなく\textbf{民間消費の直接的な減少}として現れる点に本質的な特徴があります。

\subsection{定常状態と第4回との比較}\label{ux5b9aux5e38ux72b6ux614bux3068ux7b2c4ux56deux3068ux306eux6bd4ux8f03}

外生的な技術ショックの定常水準を \(a = 0\)（すなわち
\(\exp(0)=1\)）と正規化します。定常状態において政府購入の水準を外生定数
\(G\) として与えると、定常状態の労働投入量 \(N\) は次式を満たします。 \[
N^\varphi \left( N^{1-\alpha} - G \right)^\gamma = (1-\alpha) N^{-\alpha}.
\]
この非線形方程式は、一般的なパラメータ（\(\alpha > 0, \gamma \neq 1\)）のもとでは明示的な閉形式解を得ることが困難です。

そこで、多くのマクロ経済分析の慣例に従い、定常状態における\textbf{政府支出比率
\(G_Y \equiv G/Y\)}
を所与の構造パラメータとして設定します。このとき、定常状態の民間消費は
\(C = (1-G_Y)Y\)
と極めて見通しの良い形で表せるため、解析的な定常状態解を容易に導出できます。ここで
\(G_Y\)
は定常状態を定めるための基準パラメータであり、各期において政府支出が常に産出の一定比率に固定されているという意味ではない点に注意してください。

第4回と同様の表記として \(\varphi + \alpha + \gamma(1-\alpha)\)
を分母に用いると、政府支出を含む経済の定常状態（上付き添字
\(G\)）は、以下のように一意に求まります。 \[
\begin{aligned}
N^G &= \left[ (1-\alpha) (1-G_Y)^{-\gamma} \right]^{\frac{1}{\varphi + \alpha + \gamma(1-\alpha)}},\\
Y^G &= (N^G)^{1-\alpha},\\
C^G &= (1-G_Y) Y^G,\\
W^G &= (1-\alpha) (N^G)^{-\alpha}.
\end{aligned}
\]
ここで、第4回で導出した「政府支出が存在しない経済の定常状態（上付き添字
\(0\)）」 \[
N^0 = (1-\alpha)^{\frac{1}{\varphi + \alpha + \gamma(1-\alpha)}}, \qquad Y^0 = C^0 = (N^0)^{1-\alpha}, \qquad W^0 = (1-\alpha)(N^0)^{-\alpha}
\]
と本モデルの定常状態を比較してみましょう。正の政府支出比率（\(0 < G_Y < 1\)）が存在する場合、両者の比率は次のように表されます。
\[
\begin{aligned}
\frac{N^G}{N^0} &= (1-G_Y)^{-\frac{\gamma}{\varphi + \alpha + \gamma(1-\alpha)}} > 1,\\
\frac{Y^G}{Y^0} &= (1-G_Y)^{-\frac{\gamma(1-\alpha)}{\varphi + \alpha + \gamma(1-\alpha)}} > 1,\\
\frac{C^G}{C^0} &= (1-G_Y)^{\frac{\varphi + \alpha}{\varphi + \alpha + \gamma(1-\alpha)}} < 1,\\
\frac{W^G}{W^0} &= (1-G_Y)^{\frac{\alpha\gamma}{\varphi + \alpha + \gamma(1-\alpha)}} \leq 1.
\end{aligned}
\]
この比較から、極めて明快な経済学的含意が得られます。政府が実物資源を恒常的に吸収する経済（\(G_Y > 0\)）では、家計の可処分資源が減少する負の富効果によって\textbf{労働投入量が増加（\(N^G > N^0\)）}し、その結果として\textbf{総産出水準が拡大（\(Y^G > Y^0\)）}します。一方で、生産された財の一部が政府に吸収されるため、\textbf{民間消費は確実に低下（\(C^G < C^0\)）}します。

また、労働投入の増加に伴い、収穫逓減（\(\alpha > 0\)）のもとで労働の限界生産力が低下するため、\textbf{実質賃金は低下（\(W^G < W^0\)）}します。なお、資本分配率が存在しない線形生産技術（\(\alpha = 0\)）の特殊ケースでは、労働限界生産力が労働投入量に依存しないため、実質賃金は政府支出の有無にかかわらず一定（\(W^G/W^0 = 1\)）に保たれます。

\section{対数線形近似}\label{ux5bfeux6570ux7ddaux5f62ux8fd1ux4f3c}

\subsection{対数線形均衡体系の導出}\label{ux5bfeux6570ux7ddaux5f62ux5747ux8861ux4f53ux7cfbux306eux5c0eux51fa}

経済を定常状態の近傍で対数線形近似します。まず、総資源制約式
\(Y_t = C_t + G_t\) を Taylor 展開すると、 \[
Y \frac{Y_t - Y}{Y} = C \frac{C_t - C}{C} + (G_t - G)
\] となります。両辺を定常産出 \(Y\) で割り、\(C/Y = 1 - G_Y\)
および政府支出ショックの定義 \(g_t \equiv (G_t - G)/Y\)
を代入すると、財市場の対数線形資源制約式が得られます。
\begin{equation}\protect\phantomsection\label{eq-loglinear-resource-government}{
y_t = (1-G_Y) c_t + g_t.
}\end{equation}
この式をその他の均衡条件と連立させることで、政府支出ショックを含む固定資本モデルの\textbf{対数線形体系}が以下のように確立されます。
\[
\begin{aligned}
y_t &= a_t + (1-\alpha) n_t,\\
y_t &= (1-G_Y) c_t + g_t,\\
w_t &= a_t - \alpha n_t,\\
w_t &= \gamma c_t + \varphi n_t,\\
a_t &= \rho_a a_{t-1} + e_t^a,\\
g_t &= \rho_g g_{t-1} + e_t^g.
\end{aligned}
\] 外生ショック過程において、\(e_t^a\) および \(e_t^g\)
は期待値ゼロの直交するイノベーションです。

\subsection{対数線形解の導出と解析}\label{ux5bfeux6570ux7ddaux5f62ux89e3ux306eux5c0eux51faux3068ux89e3ux6790}

連立方程式系を解き、外生ショック（\(a_t, g_t\)）に対する各内生変数の閉形式解を求めます。

まず、資源制約式から民間消費の対数偏差 \(c_t\) を表すと、 \[
c_t = \frac{y_t - g_t}{1-G_Y}
\] となります。これを労働供給条件 \(w_t = \gamma c_t + \varphi n_t\)
に代入し、さらに企業の労働需要条件 \(w_t = a_t - \alpha n_t\)
および生産関数 \(y_t = a_t + (1-\alpha)n_t\) を組み合わせて整理します。
\[
a_t - \alpha n_t = \gamma \frac{a_t + (1-\alpha)n_t - g_t}{1-G_Y} + \varphi n_t.
\] \(n_t\) の項を左辺に、\(a_t\) および \(g_t\)
の項を右辺に集めて解くと、均衡労働投入量 \(n_t\) の反応式が導かれます。
\[
n_t = \frac{(1-G_Y-\gamma) a_t + \gamma g_t}{(1-G_Y)(\alpha+\varphi) + \gamma(1-\alpha)}.
\]

後続の代表的家計モデル（RANKモデル）の標準的表記に合わせ、技術ショックに対する産出反応弾力性
\(\zeta_a\)、および政府支出ショックに対する\textbf{局所的政府購入乗数
\(\zeta_g\)} を以下のように定義します。 \[
\begin{aligned}
\zeta_a &\equiv \frac{(1-G_Y)(1+\varphi)}{(1-G_Y)(\alpha+\varphi) + \gamma(1-\alpha)} > 0,\\
\zeta_g &\equiv \frac{\gamma(1-\alpha)}{(1-G_Y)(\alpha+\varphi) + \gamma(1-\alpha)} > 0.
\end{aligned}
\] このとき、総産出 \(y_t\) は生産関数 \(y_t = a_t + (1-\alpha)n_t\)
より、極めて明快な線形結合として導出されます。 \[
y_t = \zeta_a a_t + \zeta_g g_t.
\]
価格が完全に伸縮的な本モデルにおいて決定されるこの産出水準は、第7回以降のNew
Keynesianモデルにおいて金融政策の目標基準となる「\textbf{自然産出量（Natural
Level of Output:
\(y_t^f\)）}」そのものです。また第8回では、この自然産出量に対応する消費経路
\(c_t^f\)
を家計のオイラー方程式に代入することにより、政府支出ショックを含む自然実質利子率
\(r_t^f\) を復元することになります。

求めた産出の決定式を生産関数、資源制約式、および労働需要式に代入することにより、すべてのマクロ変数の対数線形解体系が一意に確定します。
\[
\begin{aligned}
y_t &= \zeta_a a_t + \zeta_g g_t,\\
n_t &= \frac{\zeta_a - 1}{1-\alpha} a_t + \frac{\zeta_g}{1-\alpha} g_t,\\
c_t &= \frac{\zeta_a}{1-G_Y} a_t + \frac{\zeta_g - 1}{1-G_Y} g_t,\\
w_t &= \frac{1-\alpha\zeta_a}{1-\alpha} a_t - \frac{\alpha\zeta_g}{1-\alpha} g_t.
\end{aligned}
\]

\subsection{第4回（政府なしモデル）との比較と経済学的直観}\label{ux7b2c4ux56deux653fux5e9cux306aux3057ux30e2ux30c7ux30ebux3068ux306eux6bd4ux8f03ux3068ux7d4cux6e08ux5b66ux7684ux76f4ux89b3}

第4回で導出した政府支出が存在しない場合の産出反応係数を \(\zeta_a'\)
と表記すると、 \[
\zeta_a' \equiv \frac{1+\varphi}{\varphi + \alpha + \gamma(1-\alpha)}
\]
でした。本モデルにおいて定常状態の政府支出比率をゼロ（\(G_Y = 0\)）とし、政府支出ショックをゼロ（\(g_t = 0\)）と置けば、
\[
\zeta_a \Big|_{G_Y=0} = \zeta_a'
\] となり、第4回で導出したすべての解体系に厳密に一致します。

本モデルにおける最も重要な洞察は、\textbf{正の政府支出ショック（\(e_t^g > 0\)）がマクロ経済に与える影響}です。政府支出が拡大すると、

\begin{itemize}
\tightlist
\item
  総産出および労働投入量は増加します（\(y_t > 0, n_t > 0\)）。
\item
  民間消費は確実に押し下げられ、クラウディング・アウトされます（\((\zeta_g - 1)/(1-G_Y) < 0\)
  より \(c_t < 0\)）。
\item
  労働投入の拡大に伴う収穫逓減により、実質賃金は低下します（\(\alpha > 0\)
  のとき \(w_t < 0\)）。
\end{itemize}

政府支出の乖離は \(g_t \equiv (G_t - G)/Y\)
と定義され、産出の対数偏差は一次近似で \(y_t \simeq (Y_t - Y)/Y\)
と表されます。したがって、係数 \(\zeta_g\)
は、定常状態における\textbf{局所的な実質政府購入乗数（\(\partial Y_t / \partial G_t\)）}を表しています。

\Needspace{0.30\textheight}

具体的な数値例として、標準的なマクロパラメータ（\(G_Y = 0.2, \alpha = 0.4, \gamma = 1.0, \varphi = 1.0\)）のもとで乗数係数を計算してみましょう。
\[
\begin{aligned}
\zeta_g &= \frac{1.0 \times (1 - 0.4)}{(1 - 0.2)(0.4 + 1.0) + 1.0 \times (1 - 0.4)}\\[6pt]
&= \frac{0.6}{0.8 \times 1.4 + 0.6}
 = \frac{0.6}{1.12 + 0.6}\\[6pt]
&= \frac{0.6}{1.72} \simeq 0.349,\\[8pt]
\frac{\zeta_g - 1}{1-G_Y} &= \frac{0.349 - 1}{0.8} \simeq -0.814.
\end{aligned}
\] この結果は、政府支出が定常産出の \(1\%\)
分だけ増加したとき、総産出は約 \(0.349\%\)
増加する一方で、民間消費はその定常値から約 \(0.814\%\)
押し下げられることを意味します。水準値に換算すると、政府支出が \(1\)
単位増加したとき、総産出は約 \(0.349\) 単位増加し、民間消費は約
\(0.651\)
単位クラウディング・アウトされます（\(1 - 0.349 = 0.651\)）。このように乗数が
\(1\)
を大きく下回る理由は、政府支出が消費者の効用や企業の生産性を高めず、さらに投資の通時的調整が存在しないためです。資本蓄積を含む動学的一般均衡における財政乗数の詳細については、Baxter
and King (1993) や Aiyagari, Christiano, and Eichenbaum (1992)
を参照してください。

また、技術ショック \(a_t\)
に対する反応についても、政府支出の存在によって興味深い違いが現れます。\(g_t = 0\)
のもとで正の技術ショックが発生した際、対数効用（\(\gamma = 1\)）を仮定すると、第4回では代替効果と所得効果が完全に相殺して労働投入量は全く動きませんでした（\(n_t = 0\)）。しかし政府支出比率
\(G_Y > 0\) が存在する本モデルでは、 \[
n_t = -\frac{G_Y}{(1-G_Y)(\alpha+\varphi) + (1-\alpha)} a_t < 0
\]
となり、\textbf{正の技術ショックに対して労働投入量が減少する}という結果が生じます。これは、政府支出の水準が短期的に固定されているため、技術進歩による追加的な産出の伸びが民間消費へ相対的に大きく配分され、消費増加に伴う所得効果が代替効果を上回るためです。

\section{税制とラッファーカーブ}\label{ux7a0eux5236ux3068ux30e9ux30c3ux30d5ux30a1ux30fcux30abux30fcux30d6}

前章までは政府購入の財源調達を一括固定税で行うケースを考察してきました。本章では、政府購入の経路を固定した上で、限界的な意思決定に歪みをもたらす\textbf{比例税（歪みのある税制）}を導入し、税制が市場均衡に生み出す「\textbf{くさび（Tax
Wedge）}」を分析します。この分析は、第6回において独占的競争によるマークアップの歪みを売上比例補助金で相殺する政策を論じる上での不可欠な基礎となります。

\subsection{一括税と比例税の対比}\label{ux4e00ux62ecux7a0eux3068ux6bd4ux4f8bux7a0eux306eux5bfeux6bd4}

固定資本を蓄積しない本モデルにおいて比例税を導入する場合、代表的な税制として「\textbf{労働所得税}」と「\textbf{消費税}」が考えられます。本講義では、固定資本への配当収益に対する課税は除外して考えます。

まず、家計に対して労働所得税率
\(\tau_N\)（\(0 \leq \tau_N < 1\)）を課す場合を考えます。家計が労働1単位の対価として実際に受け取る税引き後実質賃金は
\((1-\tau_N)W_t\)
となるため、家計の静学的な労働供給条件は以下のように修正されます。 \[
(1-\tau_N) W_t = -\frac{U_{N,t}}{U_{C,t}}.
\] 一方で、企業の労働需要条件 \(W_t = \exp(a_t) F_N(\bar{K}, N_t)\)
は企業が支払う税込み賃金 \(W_t\)
に基づくため変化しません。その結果、労働所得税は家計の主観的限界不効用と企業の技術的限界生産力との間に
\((1-\tau_N)\) という「くさび」を打ち込みます。

同様に、家計の消費支出に対して消費税率 \(\tau_C \geq 0\)
を課す場合、財の購入価格が実質的に \((1+\tau_C)\)
倍となるため、労働供給条件は \[
\frac{1-\tau_N}{1+\tau_C} W_t = -\frac{U_{N,t}}{U_{C,t}}
\]
へと一般化されます。資本蓄積が存在しない静学的な環境においては、消費税率の引き上げ（分母の増大）と労働所得税率の引き上げ（分子の減少）は、労働供給のくさびを拡大させるという観点から全く等価な経済効果を持ちます。

\subsection{税の帰着（Tax
Incidence）の等価性}\label{ux7a0eux306eux5e30ux7740tax-incidenceux306eux7b49ux4fa1ux6027}

租税論における極めて重要な命題として、「\textbf{法律上どちらの主体に納税義務が課されているか（法的帰着）と、市場均衡において最終的に誰が税負担を負いどのような配分が実現するか（経済的帰着）は全く無関係である}」という原則があります。実物配分を一意に決定するのは、家計が受け取る限界的賃金と企業が支払う限界的労働費用との間の比率（くさびの大きさ）のみです。

これを厳密に示すため、以下の2つの課税方式を比較します。

\begin{enumerate}
\def\labelenumi{\arabic{enumi}.}
\tightlist
\item
  \textbf{家計側への労働所得課税（税率 \(\tau_N^H\)）}：
  企業が支払う実質賃金を \(W_t\) とすると、家計の税引き後受取賃金は
  \((1-\tau_N^H)W_t\) です。労働市場の均衡条件は \[
  -\frac{U_{N,t}}{U_{C,t}} = (1-\tau_N^H) \exp(a_t) F_N(\bar{K}, N_t)
  \] となります。
\item
  \textbf{企業側への労働投入課税（税率 \(\tau_N^F\)）}：
  家計が受け取る実質賃金を \(\widetilde{W}_t\)
  とすると、企業が支払う総労働費用は \((1+\tau_N^F)\widetilde{W}_t\)
  です。企業の実質利潤は \[
  D_t = \exp(a_t) F(\bar{K}, N_t) - (1+\tau_N^F) \widetilde{W}_t N_t
  \] となるため、利潤最大化の一階条件は
  \((1+\tau_N^F)\widetilde{W}_t = \exp(a_t) F_N(\bar{K}, N_t)\)
  です。家計の労働供給条件 \(\widetilde{W}_t = -U_{N,t}/U_{C,t}\)
  を代入すると、労働市場の均衡条件は \[
  -\frac{U_{N,t}}{U_{C,t}} = \frac{1}{1+\tau_N^F} \exp(a_t) F_N(\bar{K}, N_t)
  \] となります。
\end{enumerate}

これら2式の右辺を比較することにより、両者が同一の労働市場のくさびを生み出すための条件は直ちに求まります。
\begin{equation}\protect\phantomsection\label{eq-tax-incidence-equivalence}{
1 - \tau_N^H = \frac{1}{1+\tau_N^F} \qquad \Longleftrightarrow \qquad \tau_N^H = \frac{\tau_N^F}{1+\tau_N^F}.
}\end{equation} 同一の政府購入経路 \(G_t\)
が与えられ、税収が同一のルールで処理されるならば、総資源制約式
\(Y_t = C_t + G_t\)
および生産関数も同一であるため、\textbf{均衡における労働投入量
\(N_t\)、産出量 \(Y_t\)、民間消費 \(C_t\)
は完全に一致}します。法的な納税義務者が家計であるか企業であるかによって生じる違いは、名目上の賃金水準の定義の違いに過ぎません。

さらに、政府が得る総税収を比較すると、企業側課税における税収は \[
TX_t = \tau_N^F \widetilde{W}_t N_t = \frac{\tau_N^F}{1+\tau_N^F} \exp(a_t) F_N(\bar{K}, N_t) N_t = \tau_N^H \exp(a_t) F_N(\bar{K}, N_t) N_t
\]
となり、家計側課税における税収と完全に一致します。したがって、完全競争と伸縮的賃金のもとでは、労働課税の経済的帰着は完全に中立的です。

\subsection{労働所得税だけのラッファーカーブ}\label{ux52b4ux50cdux6240ux5f97ux7a0eux3060ux3051ux306eux30e9ux30c3ux30d5ux30a1ux30fcux30abux30fcux30d6}

最後に、政府購入の財源調達問題から離れ、税率と課税ベースの相互作用を明らかにするための独立した静学実験として、\textbf{労働所得税のラッファーカーブ（Laffer
Curve）}を導出します。

経済が確定的定常状態（\(a=0\)）にあり、政府債務および政府購入がともにゼロ（\(B = 0, G = 0\)）である状況を想定します。政府は家計の労働所得に対して一様な税率
\(\tau_N\)（\(0 \leq \tau_N < 1\)）を課し、得られた総税収を\textbf{全額一括移転
\(TR\) として家計に均等に還付}します。 \[
TR = TX = \tau_N W N.
\]
ここで極めて重要な仮定は、「\textbf{家計は連続体（測度ゼロの個人）から構成されており、各家計は政府から受け取る一括移転
\(TR\)
を外生的な所与のものとして行動する}」という点です。個々の家計は、自分が労働供給を増やしても国全体の総税収や移転受取額は無視できるほどしか増えないと認識するため、税収が自分自身に還付される効果を内部化しません。この「フリーライダー問題（外部性）」が存在するため、たとえマクロ全体で税収が全額還付されるとしても、労働所得税は各家計の限界的な税引き後賃金を確実に引き下げ、労働意欲を阻害します。

政府購入が存在しないため（\(G=0\)）、財市場の資源制約式は \[
C = Y
\] となります。固定資本 \(\bar{K}=1\) の生産関数 \(Y = N^{1-\alpha}\)
および企業の労働需要条件 \(W = (1-\alpha)N^{-\alpha}\)
を、労働所得税を含む家計の労働供給条件 \[
(1-\tau_N) W = N^\varphi C^\gamma
\] に代入すると、次式が得られます。 \[
(1-\tau_N) (1-\alpha) N^{-\alpha} = N^\varphi \left( N^{1-\alpha} \right)^\gamma.
\] \(N\) について整理して解くことにより、税率 \(\tau_N\)
のもとでの均衡労働投入量 \(N(\tau_N)\) が一意に求まります。 \[
N(\tau_N) = \left[ (1-\tau_N)(1-\alpha) \right]^{\frac{1}{\varphi + \alpha + \gamma(1-\alpha)}}.
\] 分母 \(\varphi + \alpha + \gamma(1-\alpha)\)
は常に正であるため、労働所得税率 \(\tau_N\)
が引き上げられると、税引き後賃金の低下を通じて均衡労働投入量
\(N(\tau_N)\) は単調に減少します。

政府の総税収は、定義より \(TX(\tau_N) = \tau_N W N\)
です。企業の分配関係 \(W N = (1-\alpha) Y = (1-\alpha) N^{1-\alpha}\)
を代入すると、税収は \[
TX(\tau_N) = \tau_N (1-\alpha) N(\tau_N)^{1-\alpha}
\] と表現されます。ここに先ほど求めた \(N(\tau_N)\)
を代入すると、労働所得税率 \(\tau_N\)
に対する税収関数、すなわち\textbf{ラッファーカーブの解析解}が導かれます。
\begin{equation}\protect\phantomsection\label{eq-laffer-revenue}{
TX(\tau_N) = \mathcal{A}\, \tau_N (1-\tau_N)^\vartheta, \qquad \vartheta \equiv \frac{1-\alpha}{\varphi+\alpha+\gamma(1-\alpha)} > 0.
}\end{equation} ここで \(\mathcal{A} \equiv (1-\alpha)^{1+\vartheta}\)
は税率に依存しない正の定数です。

この数式は、明快な逆U字型の形状を描きます。税率がゼロ（\(\tau_N = 0\)）のときは当然ながら税収はゼロです。一方で、税率が
\(100\%\)（\(\tau_N \to 1\)）に近づくと、税引き後実質賃金がゼロになるため家計は働く意欲を完全に失い、労働投入量および課税ベースがゼロに収縮して税収は再びゼロに近づきます。

税収を最大化する最適な税率 \(\tau_N^*\) は、対数税収関数の1階条件 \[
\frac{d\ln TX(\tau_N)}{d\tau_N} = \frac{1}{\tau_N} - \frac{\vartheta}{1-\tau_N} = 0
\] を解くことによって一意に導出されます。
\begin{equation}\protect\phantomsection\label{eq-laffer-peak}{
\tau_N^* = \frac{1}{1+\vartheta} = \frac{\varphi + \alpha + \gamma(1-\alpha)}{\varphi + \alpha + \gamma(1-\alpha) + (1-\alpha)}.
}\end{equation} したがって、税収の符号変化は以下のようになります。 \[
\frac{d TX}{d\tau_N} > 0 \quad (\tau_N < \tau_N^*), \qquad \frac{d TX}{d\tau_N} < 0 \quad (\tau_N > \tau_N^*).
\]
税率が低い領域（\(\tau_N < \tau_N^*\)）では、税率引き上げによる増収効果が課税ベースの縮小を上回るため税収が増加します。しかし、税率が臨界値を超えた高い領域（\(\tau_N > \tau_N^*\)）では、重税による労働意欲の減退と課税ベースの急速な縮小が税率引き上げの効果を凌駕し、\textbf{税率を引き上げるほど税収が減少する}という現象が発生します。

具体的なパラメータ値（\(\alpha = 0.4, \gamma = 1.0, \varphi = 1.0\)）のもとで計算すると、
\[
\varphi + \alpha + \gamma(1-\alpha) = 1 + 0.4 + 1 \times 0.6 = 2.0, \qquad \vartheta = \frac{0.6}{2.0} = 0.3
\] となり、理論上の最大税収率は \[
\tau_N^* = \frac{1}{1 + 0.3} = \frac{1}{1.3} \simeq 0.769 \quad (76.9\%)
\] として求まります。この \(76.9\%\)
は、資本蓄積が存在せず、労働供給の調整のみに依存する本講義の静学的モデルと、上記パラメータのもとでの値です。資本蓄積を含む
Trabandt and Uhlig (2011, Table 5) の基準モデル（Frisch労働供給弾力性
\(1\)、異時点間消費代替弾力性
\(0.5\)）では、総税収を最大にする労働所得税率は米国 \(63\%\)、EU-14
\(62\%\)
です。税収最大化税率は選好や制度設定に依存するため、本講義の数値を一般的な税率として解釈することはできません。

図~\ref{fig-laffer-curve} は、基準パラメータのもとで正規化した税収
\(TX(\tau_N)/\mathcal{A} = \tau_N(1-\tau_N)^\vartheta\)
の形状を可視化したものです。HTML版では、以下のトグルを展開してRコードを確認・実行することができます。

\begin{figure}[H]

\centering{

\pandocbounded{\includegraphics[keepaspectratio]{lecture05_files/figure-pdf/fig-laffer-curve-1.pdf}}

}

\caption{\label{fig-laffer-curve}労働所得税のラッファーカーブ（基準パラメータ：\(\alpha=0.4, \gamma=1.0, \varphi=1.0\)）。縦軸は正規化された実質税収
\(TX/\mathcal{A}\) を表し、赤い破線は最大税収率
\(\tau_N^* \simeq 0.769\)（\(76.9\%\)）を示している。}

\end{figure}%

\section{まとめと第6回への接続}\label{ux307eux3068ux3081ux3068ux7b2c6ux56deux3078ux306eux63a5ux7d9a}

本講義では、資本蓄積のない柔軟価格モデルに政府部門を導入し、財政政策の多面的な効果を分析しました。本講義の要点は以下の通りです。

\begin{enumerate}
\def\labelenumi{\arabic{enumi}.}
\tightlist
\item
  \textbf{実物資源の吸収と一括移転の峻別}：
  家計と政府の予算制約を統合すると、国債や一括税・移転などの金融的変数は相殺され、政府購入
  \(G_t\) のみが総資源制約式 \(Y_t = C_t + G_t\)
  に現れます。リカードの中立命題が成立する環境では、財源調達の選択（国債か一括税か）は実物配分に影響を与えません。
\item
  \textbf{政府支出ショックの波及効果}：
  政府支出の拡大は、民間が利用可能な資源を奪う負の富効果をもたらし、労働供給を促進して総産出を増加させる一方で、民間消費を直接的にクラウディング・アウトします。
\item
  \textbf{税制の歪みとラッファーカーブ}：
  比例的な労働課税は労働市場の限界条件に「くさび」を打ち込み、法的帰着にかかわらず同一の歪みをもたらします。税率引き上げは税率効果と課税ベース縮小効果のトレードオフを内包しており、税収と税率の間には逆U字型のラッファーカーブが成立します。
\end{enumerate}

\subsection{第6回への接続：独占的競争と最適補助金政策}\label{ux7b2c6ux56deux3078ux306eux63a5ux7d9aux72ecux5360ux7684ux7af6ux4e89ux3068ux6700ux9069ux88dcux52a9ux91d1ux653fux7b56}

第5回までのモデルでは、財市場および労働市場が完全競争であると仮定してきました。次回の第6回講義では、経済に\textbf{差別化された中間財を生産する独占的競争企業}を導入します。独占力を持つ企業は価格を限界費用以上に設定するため、市場には「\textbf{価格マークアップのくさび}」が生じ、産出量と雇用が社会的に望ましい水準を下回るという市場の失敗が発生します。

この独占的競争の歪みを是正するために用いられるのが、本講義で学んだ税・財政政策の応用です。政府は中間財企業に対して\textbf{売上比例補助金}を支給することで企業の限界費用を引き下げ、マークアップのくさびを相殺して効率的な生産を回復させます。その補助金の財源は一括税によって調達されます。一括税は限界的意思決定を歪めないため、本講義の資源制約式
\(Y_t = C_t\)（実質資源の吸収がない純粋な所得再配分）を保ったまま、生産側の歪みだけを完璧に治療することができるのです。

\clearpage

\section{参考文献}\label{ux53c2ux8003ux6587ux732e}

予習は講義ノートを中心に進めてください。以下の文献は、各項目に記した論点を復習・補足するための参照先です。

\textbf{基本参考文献}

\begin{itemize}
\tightlist
\item
  Romer, David (2019) \emph{Advanced Macroeconomics} (5th ed.).
  McGraw-Hill.
  \textbf{政府支出ショック、租税の経済効果、およびリカード等価定理の最も標準的な教科書解説}
\item
  Barro, Robert J. (1974) ``Are Government Bonds Net Wealth?''
  \emph{Journal of Political Economy}, 82(6), pp.~1095-1117.
  \href{https://doi.org/10.1086/260266}{doi}.
  \textbf{リカードの中立命題を現代マクロ経済学に定着させた記念碑的論文}
\item
  Galí, Jordi (2015) \emph{Monetary Policy, Inflation, and the Business
  Cycle: An Introduction to the New Keynesian Framework and Its
  Applications} (2nd ed.). Princeton University Press.
  \textbf{第3章が基本NKモデルの標準的参照先であり、政府支出を導入する第8回への橋渡し}
\end{itemize}

\textbf{発展文献}

\begin{itemize}
\tightlist
\item
  Baxter, Marianne, and Robert G. King (1993) ``Fiscal Policy in General
  Equilibrium.'' \emph{American Economic Review}, 83(3), pp.~315-334.
  \href{https://people.bu.edu/mbaxter/papers/fpol.pdf}{著者公開版}.
  \textbf{資本蓄積を内生化した動学的一般均衡における財政乗数の先駆的研究}
\item
  Aiyagari, S. Rao, Lawrence J. Christiano, and Martin Eichenbaum (1992)
  ``The Output, Employment, and Interest Rate Effects of Government
  Consumption.'' \emph{Journal of Monetary Economics}, 30(1), pp.~73-86.
  \href{https://doi.org/10.1016/0304-3932(92)90045-4}{doi}
\item
  Leeper, Eric M. (1991) ``Equilibria under `Active' and `Passive'
  Monetary and Fiscal Policies.'' \emph{Journal of Monetary Economics},
  27(1), pp.~129-147.
  \href{https://doi.org/10.1016/0304-3932(91)90007-B}{doi}.
  \textbf{物価水準の財政理論（FTPL）の基礎となる金融・財政レジームの分類}
\item
  Trabandt, Mathias, and Harald Uhlig (2011) ``The Laffer Curve
  Revisited.'' \emph{Journal of Monetary Economics}, 58(4), pp.~305-327.
  \href{https://doi.org/10.1016/j.jmoneco.2011.07.003}{doi}.
  \textbf{資本蓄積と内生的成長を考慮したラッファーカーブの定量的推計}
\end{itemize}

\clearpage

\section{演習問題}\label{ux6f14ux7fd2ux554fux984c}

\subsection{基本問題}\label{ux57faux672cux554fux984c}

\textbf{問1
［Core］［Derivation］：政府の通時的予算制約の導出（確実性下の基礎）}

確実性の世界において、実質総利子率が一定値 \(R^f > 1\)
で推移し、政府債務が公債動学方程式 \[
B_{t+1} = R^f B_t - PB_{t+1}
\] に従うとします。公債に関する横断性条件（TVC） \[
\lim_{i \to \infty} \frac{B_{t+i}}{(R^f)^i} = 0
\] のもとで前向きに逐次反復代入を行い、現在の債務水準が \[
B_t = \sum_{i=1}^\infty \frac{PB_{t+i}}{(R^f)^i}
\] として表されることを厳密に導出しなさい。また、右辺の級数の和が \(t\)
期ではなく \(t+1\)
期から始まる経済学的理由を述べなさい。さらに、不確実性が存在する一般的な経済において、1期先のオイラー方程式
\(B_t = \mathbb{E}_t [Q_{t,t+1}(B_{t+1} + PB_{t+1})]\)
を反復適用することで、本文の通時的予算制約式（式~\ref{eq-government-intertemporal-budget}）が得られる論理ステップを説明しなさい。

\textbf{問2 ［Core］：一括税と政府購入の経済効果の相違}

本講義の資本蓄積のない固定資本モデルにおいて、一括税・一括移転と実質政府購入
\(G_t\)
が実物配分（産出、消費、労働、実質賃金）に与える影響が本質的に異なる理由を、資源制約式と家計の一階条件に触れて平易に説明しなさい。また、国債発行のタイミングが実物配分に対して中立となるための前提条件を3つ以上挙げ、それらの仮定が現実経済において破られる具体的な状況を1つ示しなさい。

\Needspace{0.25\textheight}

\textbf{問3 ［Core］［Derivation］：対数線形資源制約式の導出}

財市場清算条件（総資源制約式）\(Y_t = C_t + G_t\)（式~\ref{eq-resource-constraint-government}）および定常産出で規格化した政府支出ショックの定義式
\(g_t \equiv (G_t - G)/Y\)（式~\ref{eq-government-spending-shock}）から、定常状態近傍での1次Taylor展開を用いて、対数線形化された資源制約式
\[
y_t = (1-G_Y) c_t + g_t
\]
（式~\ref{eq-loglinear-resource-government}）を途中の計算過程を含めて導出しなさい。

\subsection{発展問題}\label{ux767aux5c55ux554fux984c}

\textbf{問4 ［Research
extension］［Computation］：政府購入乗数とクラウディング・アウト}

正の一時的政府支出ショック（\(g_t > 0\)）が発生したとき、民間消費、労働投入量、総産出量、および実質賃金がそれぞれどのように変動するかを、代替効果と所得効果の観点から説明しなさい。また、本モデルにおいて「クラウディング・アウト」が投資ではなく民間消費に現れる理由を述べなさい。さらに、パラメータ値が
\(G_Y = 0.2, \alpha = 0.4, \gamma = 1.0, \varphi = 1.0\)
であるとき、実質政府購入乗数 \(\zeta_g\) および民間消費の反応係数
\((\zeta_g - 1)/(1-G_Y)\)
の数値を計算し、その結果の経済学的意味を解説しなさい。

\textbf{問5 ［Research
extension］［Derivation］：労働課税における税の帰着の等価性}

家計側に課される労働所得税率 \(\tau_N^H\) と企業側に課される労働投入税率
\(\tau_N^F\)
について、両者が同一の実物配分（労働投入、産出、消費）をもたらすための等価条件
\[
1 - \tau_N^H = \frac{1}{1+\tau_N^F}
\]
（式~\ref{eq-tax-incidence-equivalence}）を導出しなさい。また、この条件が満たされているとき、政府が得る総税収も両者のもとで完全に一致することを示しなさい。

\textbf{問6 ［Research
extension］［Derivation］：ラッファーカーブと最大税収率の導出}

政府購入と政府債務がともにゼロ（\(G = 0, B = 0\)）であり、労働所得税収が全額一括移転として家計に還付される（\(TR = TX\)）定常状態経済を考えます。各家計が一括移転を所与として行動する理由を述べた上で、均衡労働投入量
\(N(\tau_N)\)、総税収関数
\(TX(\tau_N)\)（式~\ref{eq-laffer-revenue}）、および最大税収率
\(\tau_N^*\)（式~\ref{eq-laffer-peak}）を順次導出しなさい。さらに、税率が
\(\tau_N^*\)
を超える高い領域において、税率を引き上げると逆に税収が減少する直観的メカニズムを説明しなさい。

\Needspace{12\baselineskip}

\textbf{問7 ［Research extension］：第6回モデルへの接続と売上比例補助金}

次回の第6回講義において、独占的競争企業の価格マークアップを是正するために売上比例補助金を導入し、その財源を一括税で賄う場合を考えます。このとき、本講義の政府購入
\(G_t\) の場合とは異なり、なぜ総資源制約式に政府の支出項目が現れず
\(Y_t = C_t\)
が維持されるのかを説明しなさい。また、売上比例補助金が企業のどの限界条件にどのような「くさび」を打ち込み、独占の歪みを相殺するのかを論じなさい。




\end{document}
